\begin{question}

A certain observable in quantum mechanics has a $3 \times 3$ matrix
representation as follows:

\begin{equation*}
  \frac{1}{\sqrt{2}}
  \begin{pmatrix}
    0 & 1 & 0\\
    1 & 0 & 1\\
    0 & 1 & 0\\
  \end{pmatrix}
\end{equation*}

\begin{questionpart}
  Find the normalized eigenvectors of this observable and the
  corresponding eigenvalues. Is there any degeneracy?
\end{questionpart}

\begin{questionpart}
 Give a physical example where all this is relevant.  
\end{questionpart}

\end{question}

\begin{purpose}
\end{purpose}

\begin{answer}

  \begin{answerpart}{Interesting Eigenvalues \& Rotations}

    Let's start by getting rid of the constant since it's just
    annoying.  Since we can always pull out a scalar ($c_0$ in this
    case), we're going to define our matrix without it and we can
    apply it at the end.

    \begin{equation}
      \text{let } A \equiv
      \begin{pmatrix}
        0 & 1 & 0\\
        1 & 0 & 1\\
        0 & 1 & 0\\
      \end{pmatrix}
    \end{equation}

    \begin{equation}
      \text{let } c_0 \equiv \frac{1}{\sqrt{2}}
    \end{equation}

    Obviously, this matrix is going to do some interesting stuff, so
    let's start with a few observations about the symmetry of the
    matrix:

    \begin{equation}
      A = A^\dagger = A^* = A^T
    \end{equation}

    We can also rotate the matrix like a pinwheel about the cetner 0
    and we get the same matrix.  So we have bilateral and radial
    symmetry in our matrix.

    If we start taking the matrix to various powers, we see something
    interesting happen.

    \begin{equation}
      \begin{split}
        A^1 &=
        \begin{pmatrix}
          0 & 1 & 0\\
          1 & 0 & 1\\
          0 & 1 & 0\\
        \end{pmatrix}\\
        A^2 &=
        \begin{pmatrix}
          1 & 0 & 1\\
          0 & 2 & 0\\
          1 & 0 & 1\\
        \end{pmatrix}\\
        A^3 &=
        \begin{pmatrix}
          0 & 2 & 0\\
          2 & 0 & 2\\
          0 & 2 & 0\\
        \end{pmatrix}\\
        A^4 &=
        \begin{pmatrix}
          2 & 0 & 2\\
          0 & 4 & 0\\
          2 & 0 & 2\\
        \end{pmatrix}\\
      \end{split}
    \end{equation}

    If $n$ is odd, we find that:

    \begin{equation}
      A^n = 2^{\frac{n-1}{2}}
      \begin{pmatrix}
        0 & 1 & 0\\
        1 & 0 & 1\\
        0 & 1 & 0\\
      \end{pmatrix}
    \end{equation}

    and if $n$ is even, we get

    \begin{equation}
      A^n = 2^{\frac{n}{2} - 1}
      \begin{pmatrix}
        1 & 0 & 1\\
        0 & 2 & 0\\
        1 & 0 & 1\\
      \end{pmatrix}\\
    \end{equation}

    You can trace out how this happens if you look closely, but it
    kinda feels like we're in the classic CS game \emph{The Game of
    Life}.

    Moving on to the Eigenstage, we can easily find the eigenvalues
    using the classical approach:

    \begin{equation}
      \begin{split}
        0 &=
        det \left(
        \frac{1}{\sqrt{2}}
        \begin{bmatrix}
          0-\lambda & 1 & 0\\
          1 & 0-\lambda & 1\\
          0 & 1 & 0-\lambda\\
        \end{bmatrix}
        \right)\\
        &= -\lambda^3 - (-\lambda - \lambda)\\
        &= -\lambda^3 + 2\lambda\\
        &= (\lambda - 0)(\lambda - \sqrt{2})(\lambda + \sqrt{2})\\
      \end{split}
    \end{equation}

    It's worth noting that we have $(-\sqrt{2}, 0, \sqrt{2})$ as our
    Eigenvalues.  If we were to apply our $c_0$ from the beginning,
    then those Eigenvalues become $(-1, 0, 1)$ which starts to look
    physically interesting.  It's still just a constant, though.  Note
    that we have both an Eigenvalue of $0$ and also a singular matrix
    $A$.  Now let's solve for our (non-normalized but easy-to-read)
    eigenvectors.

    \begin{equation}
      \begin{split}
        A \vec{x} &= \lambda\vec{x}\\
        \begin{pmatrix}
          0 & 1 & 0\\
          1 & 0 & 1\\
          0 & 1 & 0\\
        \end{pmatrix}
        \begin{pmatrix}
          x\\
          y\\
          z\\
        \end{pmatrix}
        &=
        \lambda
        \begin{pmatrix}
          x\\
          y\\
          z\\
        \end{pmatrix}\\
      \end{split}
    \end{equation}

    From this, we can easily build our characteristic equations:

    \begin{equation}
      \inlabel{eigEquations}
      \begin{split}
        y &= \lambda x\\
        x + z &= \lambda y\\
        y &= \lambda z\\
      \end{split}
    \end{equation}

    Using the first and last, we can easily see that $x = z$.
    Defining $x/z$ to be $1$ for simplicity, we can rewrite our
    equations:

    \begin{equation}
      \begin{split}
        x &= 1\\
        2 &= \lambda y => y = \frac{2}{\lambda}\\
        z &= 1\\
      \end{split}
    \end{equation}

    Knowing, as we do, that our non-degenerate $\lambda's$ are $\pm
    \sqrt{2}$, we can go ahead and rationalize this to be:

    \begin{equation}
      \begin{split}
        x &= 1\\
        y &= \pm \frac{2}{\sqrt{2}} = \pm\sqrt{2} = \lambda\\
        z &= 1\\
      \end{split}
    \end{equation}

    In the end, our Eigenvectors become rather interesting:

    \begin{equation}
      \vec{a_{1,3}} = 
      \begin{pmatrix}
        1\\
        \lambda\\
        1\\
      \end{pmatrix}\\
    \end{equation}

    Returning back to our $0$ Eigenvalue and using equation
    \inref{eigEquations}, we see that:

    \begin{equation}
      \begin{split}
        y &= 0\\
        x + z &= 0 => x = -z\\
        y &= 0\\
      \end{split}
    \end{equation}

    Once again, we'll define $x = 1$ for simplicity, which makes our
    final Eigenvector
    
    \begin{equation}
      \vec{a_2} = 
      \begin{pmatrix}
        1\\
        \lambda\\
        -1\\
      \end{pmatrix}\\
    \end{equation}

    It's worth noting that except for that pesky negative sign on the
    z component, we almost got away with a single expression for the
    Eigenvectors.

    Since the book asks us to normalize, let's do that and also bring
    back our $c_0$ as part of it.  So we're going to normalize our
    eigenvector to have magnitude  trivially adding a
    constant out front

    \begin{equation}
      c = \frac{1}{c_0\abs{a}} = \frac{1}{2c_0} = \frac{\sqrt{2}}{2}
    \end{equation}

    This constant ensures that the Eigenvectors we found are
    normalized for the scaled matrix $c_0A$ specified in the problem.
    Pull out the $\frac{1}{c_0}$ or leave it in depending upon whether
    you want the vectors from before or after (respectively) the $c_0$
    is applied.

    \begin{equation}
      \vec{a_{1,3}} =
      \frac{\sqrt{2}}{2}
      \begin{pmatrix}
        1\\
        \lambda\\
        1\\
      \end{pmatrix}\\
    \end{equation}

    Then for our final Eigenvector for our Eigenvalue of $0$

    \begin{equation}
      \vec{a_{2}} =
      \frac{\sqrt{2}}{2}
      \begin{pmatrix}
        1\\
        \lambda\\
        -1\\
      \end{pmatrix}\\
    \end{equation}
  \end{answerpart}
  
  \begin{answerpart}{Physical Interpretation}
    There are a few cases where I can intuit this being useful, such
    as with raising and lowering.  There's also clearly something to
    the scaled 2-state aspect where $A^{n+2} = 2A^{n}$.  That being
    said...I don't see any specific case from chapter 1 where I could
    use this thing.

    Employing external help, it turns out that this is a spin operator
    for a Spin-1 system which is well beyond the scope of Chapter 1.
    It's at times like this students such as myself lament Sakurai's
    premature passing.  Since this portion is so far beyond the scope
    of chapter 1, I'm going to leave it alone at that.
  \end{answerpart}

\end{answer}
